% v3-basic.tex -- V1-re gate PoC 1 (plain article, no jlreq)
% bxcjkvert vertical body text + the two bxcjkvert-specific behavioral
% fixes vs raw CJKvert: baselinestretch retention and no \normalfont
% reset on direction switch. Multi-paragraph flow + savepos probing of
% glyph/column advance direction.
\documentclass[a5paper]{article}
\usepackage[margin=20mm]{geometry}
\usepackage[utf8]{inputenc}
\makeatletter\def\@begindocumenthook{\CJKvert}\makeatother
\usepackage[ipaex-type1]{bxcjkjatype}
\usepackage[main=retain]{bxcjkvert}
\makeatletter
\begin{document}
\begin{CJK}{UTF8}{min}
\makeatletter
\typeout{V3-BLSTRETCH-BEFORE: \baselinestretch}
\typeout{V3-FONT-BEFORE: \f@family/\f@series/\f@shape, \f@size pt}
\CJKvert
\typeout{V3-DIR-AFTER-VERT: \CJK@direction}
\typeout{V3-BLSTRETCH-AFTER: \baselinestretch}
\typeout{V3-FONT-AFTER: \f@family/\f@series/\f@shape, \f@size pt}
私はその人を常に先生と呼んでいた。だからここでもただ先生と書くだけで
本名は打ち明けないでしまう。これは世間を憚かる遠慮と言うより、
その方が私にとって自然だからである。\par
私はその人の記憶を呼び起すごとに、すぐ「先生」と言いたくなる。
筆を執っても心持は同じ事である。よそよそしい頭文字などはとても
使う気にならない。\par
\begingroup
\pdfsavepos\write\@auxout{\string\gdef\string\vposA{\the\pdflastxpos,\the\pdflastypos}}%
阿\pdfsavepos\write\@auxout{\string\gdef\string\vposB{\the\pdflastxpos,\the\pdflastypos}}%
阿\pdfsavepos\write\@auxout{\string\gdef\string\vposC{\the\pdflastxpos,\the\pdflastypos}}\par
\endgroup
\CJKhorz
\typeout{V3-DIR-BACK: \CJK@direction, \baselinestretch}
水平方向に戻った。ABC 123。
\end{CJK}
\end{document}
